% Generated by task tex-groundtruth from dossiers/gt-table-pub.md (section 8, with the parts of
% sections 4 to 6 its findings cite), with the corrections of dossiers/verify-gt-table-pub.md.
\section{The table sumcheck and the public input: findings}\label{sec:gen-table-pub-findings}

Short forms used in this section. The number after the colon is the line; leanVM files are at
the pin \code{a386121f}, leanerVM files at \code{b435631}; a bare \code{:N} is a line of the file
last named.

{\small\begin{tabular}{@{}ll@{}}
\code{bp:} & \file{docs/roadmap/protocol-blueprint.md} (leanerVM) \\
\code{Seams.lean}, \code{Compose.lean} & \file{LeanerVM/Protocol/Spine/Seams.lean}, \file{…/Spine/Compose.lean} (leanerVM) \\
\code{PublicInput.lean} & \file{LeanerVM/Protocol/PublicInput.lean} (leanerVM) \\
\code{03:}, \code{05:}, \code{07:}, \code{08:} & \file{doc/leanvm/body/03-proving-primitives.tex}, \file{05-arithmetization.tex}, \\
 & \file{07-instruction-tables.tex}, \file{08-end-to-end-protocol.tex} \\
Rust files & under \file{crates/lean_vm/src/}, except the following \\
\code{transcript.rs} & \file{crates/fiat_shamir/src/transcript.rs} \\
\code{primitives/src/…} & \file{crates/primitives/src/…} \\
\code{py:} & \file{python-verifier/verifier.py} \\
\code{aggregate.py} & \file{crates/rec_aggregation/guests/aggregate.py}, the recursion guest \\
\end{tabular}}

Blueprint lines are at \code{b435631}. (The citation check noted that at the merged
\code{HEAD}, \code{8d3ea7d}, every blueprint line cited here is 4 higher: the blueprint gained a
four-line paragraph after line 51 and a rewritten conventions row \emph{Module system}.) The
step numbers (T0 to T12, P0 to P4) refer to
\cref{tab:gen-table-pub-transcript-table,tab:gen-table-pub-transcript-pub}.

%------------------------------------------------------------------------------------------------
\begin{finding}{major}{f:table-pub-bus-seam}{The bus seam admits statements the deployed table sumcheck cannot serve}
\fhead{Evidence} The dossier's section 6 (E.1), its probes (section 9) and ArkLib's definition of
round-by-round knowledge soundness (\code{rbrKnowledgeSoundnessWorstCaseWith},
\file{RoundByRound.lean:553-568}, ArkLib \code{dca90385}), which quantifies over every input
statement. The seam as built:

\begin{leancode}
def bus := of I fun (s : I.Stmt × BusOut I) q ↦ (∀ c ∈ s.2.linear, c.Holds q) ∧
  (∀ c ∈ s.2.linear, ∀ t ∈ c.terms, t.poly.totalDegree ≤ I.d) ∧
  (∀ c ∈ s.2.columns, c.Holds q) ∧ I.PublicLinesHold s.1 q ∧ I.aux q
\end{leancode}
\src{LeanerVM/Protocol/Spine/Seams.lean:175-177 (b435631)}

The consequence the dossier uses: if for some statement outside the input seam a prover makes the
verifier accept with an output inside the output seam with probability $p$, then
$Σ_i \mathit{err}\ i ≥ p$. The deployed verifier takes one point $ζ$, one form per side and
table, three values, and reads the constraints off the instance. \lean{Seam.bus} admits four
things it does not handle.
\begin{enumerate}[label=(\alph*)]
\item \emph{A clause on the statement alone, which the verifier must then check.} The degree
  clause does not mention $q$. Probe: on the toy, the stack \code{q0} with columns \code{[1,1]},
  \code{[1,1]}, \code{[0,0]} satisfies \code{M3Holds toy 0}; the statement with the single claim
  ``the term $X₂³$ at the point 0 sums to 0'' is true of \code{q0}, its public line holds, and it
  is outside \code{Seam.bus toy} by the degree clause alone. On that statement the summand is
  identically zero: a verifier that does not inspect degrees accepts the all-zero run with true
  column claims, with probability 1. So a table verifier must compute \lean{totalDegree} of every
  term it is handed and reject above \lean{I.d}. The deployed verifiers have no such check and no
  such input. In the composed protocol the guard never fires.
\item \emph{Any number of linear claims: the error on $ξ$ is unbounded over the seam.} Take $k$
  claims, each one term with the constant polynomial 1 (degree 0), weight 1: true sum 1. Give
  claim $c$ the value $1 + δ_c$ with $Σ δ_c X^c = ∏_{i<k−1}(X + a_i)$ for distinct
  $a_i ∈ \E$. The statement is outside the seam, and at each of the $k − 1$ roots the batched
  target is the true batched sum, after which the honest prover wins. So
  $\mathit{err}\ ξ ≥ (k−1)/|\E|$ for every $k$ the verifier accepts. With
  $\mathit{err}\ ξ = (B+2)/|\E|$ the verifier must reject more than $B + 3$ claims. The probe
  exhibits a list of 1000 true claims on the toy, which has one constraint and one flush.
\item \emph{Terms of one table at unrelated points.} The probe's claim \code{twoPoints} has two
  terms of table 0, at the points 0 and 1, and is inside the seam (each clause evaluated). The
  deployed verifier computes one factor $∏(1+ζ_m+r_m)$ per table. Completeness on
  \code{twoPoints} needs one factor per term.
\item \emph{Terms on any table of the instance.} With Layer 3's instance
  (\cref{f:table-pub-sumcheck-tables}), a term may sit on the memory's table, of log-height
  $κ_{\mathrm{mem}} ≥ 16$. A sumcheck of $τ_{\max}$ rounds, $τ_{\max}$ the maximum over the opcode
  tables, cannot fold it when $τ_{\max} < κ_{\mathrm{mem}}$.
\end{enumerate}
(The citation check could not re-run the two Lean probes: no Lean may be run in the merged
checkout. It found the probe sources reproduced in the dossier byte-identical to
\file{probes/gt-table-pub/SeamBusShape.lean} and \file{SeamBusMember.lean}; they use only the
numerals 0 and 1 in $\K$, so the change of meaning of numerals at the new CompPoly pin does not
affect them; whether they still elaborate on the new pins is unknown.)

\fhead{Classification} An error of the blueprint. The seam is more general than any phase with
leanVM's transcript and leanVM's checks can serve. \lean{Phase.Complete} and
\lean{Phase.Security} for the deployed verifier against \lean{Seam.bus} as built cannot be proved:
they need guards (degree, number of claims, one point per table, tables of the sumcheck) that the
deployed verifier has not.

\fhead{Proposed change}
\begin{change}{the bus seam takes leanVM's shape}
\asis (\code{bp:481}, \code{bp:487-489}, \code{bp:560-563}, \code{bp:565-574};
\code{Seams.lean:71-95, 130-134, 175-177}) \lean{BusOut I} is
\lean{linear : List (LinearClaim I)}, \lean{columns}; ``The seams carry \emph{claims}, not
challenges … which is what lets the table sumcheck consume any list of linear claims without
knowing the bus.''
\asproposed The statement has leanVM's shape, the Rust's \code{BusVerify}
(\cref{f:table-pub-layer7-sketch}).
\begin{leancode}
/-- A polynomial of the row of table `j` within the degree bound. -/
abbrev M3Instance.RowPoly (I) (j) := {p : CMvPolynomial (I.width j) K // p.totalDegree ≤ I.d}
/-- A bus form of table `j`: a combination of row polynomials with weights in `E`. -/
abbrev M3Instance.Form (I) (j) := List (E × I.RowPoly j)
structure BusOut (I : M3Instance) where
  point   : Vector E I.τmax                       -- ζ_{<τ_max}
  forms   : Fin 3 → (j : I.SumcheckTable) → I.Form j   -- push, pull, count
  totals  : Fin 3 → E
  columns : List (ColumnClaim I)
\end{leancode}
\lean{Seam.bus}: every constraint of every sumcheck table has extension zero at
$\mathit{point}_{<τ_j}$; for every side, the forms' extensions at $\mathit{point}_{<τ_j}$ sum to
\lean{totals}; every column claim holds; the public lines; \lean{aux}. The prose becomes: ``The
bus seam carries the point the bus phase ended at, since the table sumcheck runs at it (§5.5);
$α, β$ and the leaf layout stay inside the bus phase, which hands on the forms it built.''
\reason The phase is then leanVM's, with error $(B+2)/|\E|$ for $B$ read off $I$, and its
verifier has no check leanVM has not. The degree bound moves from a clause to a type.
\end{change}
\emph{Alternative, smaller}: keep the types and add to \lean{Seam.bus} the clauses ``at most
$B + 3$ linear claims'', ``the terms of a table share their point and the points are prefixes of
one vector'', ``terms sit on sumcheck tables''; state in Layer 7 that the verifier checks them and
the degree, and in Layer 12 that the four guards are dead in the composition.
\end{finding}

%------------------------------------------------------------------------------------------------
\begin{finding}{major}{f:table-pub-sumcheck-tables}{The table sumcheck's tables are not distinguished from the instance's other tables}
\fhead{Evidence} Blueprint \code{bp:836-839}, ``The stack columns that belong to no opcode table
(\code{mem_0, mem_1, mem_2}, \code{cntfin_mem}, \code{cntfin_bc}, \code{q_flock} …) are tables
of the instance with no constraints and no flushes''; Layer 7 \code{bp:987-989}, $τ_{\max}$ and
$Σ_j \mathit{width}_j$ over the instance; ground truth T4 and T9. With $κ_{\mathrm{mem}} ≥ 16$
every leanISA instance would have at least 16 rounds and 110 final values; leanVM has $τ_{\max}$
rounds (as few as 3) and 104. \textbf{Unverified}: \lean{leanIsaInstance} is not built; this is a
reading of the blueprint's text.

\fhead{Classification} An error of the blueprint (an omission).

\fhead{Proposed change}
\begin{change}{say which tables the sumcheck ranges over}
\asis (row \emph{Table sumcheck}, \code{bp:323}) ``Variables bound highest first; table $j$
joins at round $τ_{\max} − τ_j$; …''
\asproposed Prepend ``The sumcheck ranges over the tables that have a constraint, a flush or a
count column (\lean{I.SumcheckTable}; for leanISA the six opcode tables, never the tables of the
shared columns); $τ_{\max}$ is the largest log-height among them and the final message has one
value per column of each, in table order then column order.'' Add \lean{SumcheckTable} (a
decidable predicate on \lean{Fin I.ntab}, derived from \lean{constraints}, \lean{flushes},
\lean{counts}) to the spine, and an acceptance test: on an instance with a taller table of no
constraint and no flush, the number of rounds is the smaller height.
\reason The schedule is Category B; a longer one rejects every Rust proof.
\end{change}
\end{finding}

%------------------------------------------------------------------------------------------------
\begin{finding}{major}{f:table-pub-public-input-check}{The public-input phase proves a check no executable verifier runs, and the debt is assigned to a layer that cannot pay it}
\fhead{Evidence} The ground truth of the public-input check
(\cref{sec:gen-table-pub-transcript-pubcheck}). New to the record: the recursion guest has the
combined check too (\code{aggregate.py:1683}).

\fhead{Classification} A deliberate deviation (decision 15; the status file's finding on the
combined public-input check (F18)). It owes a theorem: round-by-round knowledge soundness of the
verifier that checks the combined equation and pools the scalars sent. The blueprint says
``Layer 12 owes it'' (\code{bp:1062}). Layer 12 is the compilation. The combined check is a
different \emph{oracle verifier}; \lean{verify_knowledgeSound} is derived from the oracle
protocol \lean{verify} compiles; so either \lean{verify} has the per-limb check and is not the
deployed verifier (it rejects proofs the Rust accepts), or it has the combined check and the
oracle protocol must have it, with its \lean{Phase.Security}.

\fhead{Proposed change}
\begin{change}{a second public-input phase for the deployed check}
\asis (\code{bp:1060-1062}) ``So the pinned verifiers accept transcripts this verifier rejects,
their knowledge soundness at $1/|\E|$ is a lemma of its own, and Layer 12 owes it.''
\asproposed ``So the pinned verifiers accept transcripts this verifier rejects. Layer 8 has a
second phase for them, \lean{publicInputPhaseWords}: the same schedule; the check
$Σ_ℓ y^ℓ·c_ℓ = Σ_ℓ y^ℓ·\mathit{lineValue}_ℓ(r)$ over the lines whose value is sent, at most
three; the claims pooled at the values \emph{sent}. Its \lean{Phase.Security} from
\lean{Seam.table} to \lean{Seam.pub} has error $1/|\E|$: the state after the challenge is `the
true evaluations pass the check and every pooled claim holds', and a wrong memory gives
$A + r·(A + C) = 0$ with $(A, C) ≠ (0, 0)$ by the independence of $1, y, y²$ over $\K$. The
leanISA protocol and \lean{verify} use this phase; the per-limb phase is kept as the
specification's.''
\reason Three verifiers against one paragraph; the master theorem must be about the verifier that
is deployed. The discrepancy is also a finding against the specification, to report to leanVM
when reporting is allowed.
\end{change}
\end{finding}

%------------------------------------------------------------------------------------------------
\begin{finding}{minor}{f:table-pub-round-message}{The round message: four coefficients and a check in the oracle protocol, three and none on the wire}
\fhead{Evidence} Steps T5 and T6; blueprint \code{bp:315}.

\fhead{Classification} A deliberate deviation. It owes a lemma the blueprint does not state.
``Inverted by the running claim'' holds only on messages that pass the round check; and ArkLib's
\lean{Verifier.fiatShamir} of a protocol with four-coefficient messages derives the challenge from
four coefficients, where leanVM's chain absorbs three (\code{transcript.rs:303-307}). So
\lean{verify_iff_compiled} as sketched (\code{bp:1218-1221}) needs a challenge oracle defined on
the \emph{encoded} message, and \lean{FiatShamirSecurity} must be stated for that oracle.

\fhead{Proposed change} One of:
\begin{itemize}
\item make the wire's message the oracle protocol's (three coefficients, the verifier derives the
  fourth, no round check): the compilation is then literal and one check leaves the audit;
\item or keep row \emph{Sumcheck messages} and add: ``Layer 4 proves the transport: if a verifier
  $V$ reading messages $m$ is round-by-round knowledge sound and $\mathit{dec}$ is injective from
  wire messages onto the messages passing $V$'s round check, commuting with the challenges, then
  $V ∘ \mathit{dec}$ is, at the same error, with the state function
  $\mathit{state} ∘ \mathit{dec}$.''
\end{itemize}
\end{finding}

%------------------------------------------------------------------------------------------------
\begin{finding}{minor}{f:table-pub-zerocheck-escape}{The zerocheck escape is over-charged, and charged in three different ways}
\fhead{Evidence} What is stated (dossier \file{gt-table-pub.md}, section 5, D.1). The specification
gives no error for the recycled point: \code{05:132} says it is sound ``as long as it results from
uniform sampling occurring after the table columns are committed''. The blueprint charges it in
three ways: row \emph{Seams} (\code{bp:333}): ``$1/|\E|$ per coordinate per constraint'', in the
bus phase; Layer 7 (\code{bp:1001-1003}): ``charged to the $(α, β)$ and GKR challenges of Layer 6
… as an extra $τ_{\max}/|\E|$ per constraint''; Layer 6 (\code{bp:973-975}): \lean{busError} ``is
$4·2^{μ_{\mathrm{bus}}}/|\E|$ on the $(α, β)$ message plus \lean{gkrError 3 μ_bus}'', no such
term. Acceptance test 6 (\code{bp:1283}): ``The state function of Layer 7 charges the zerocheck to
Layer 6's challenges''. The tracker's hole comment for the bus phase repeats ``$1/|\E|$ per
coordinate per constraint''.

The recycled point, settled (dossier section 5, D.3). The last GKR layer is radix 4
($μ_{\mathrm{bus}} ≥ 16$); $ζ = (\mathit{low}, \mathit{high}, x_0, …, x_{μ_{\mathrm{bus}}−3})$,
so $ζ_0, ζ_1$ are the last coordinates of $ζ$ drawn and $ζ_{2+i} = x_i$. (The citation check
corrected the dossier's ``are drawn last'': after $(\mathit{low}, \mathit{high})$ the verifier
draws one more combiner, unused after the last layer, \code{gkr.rs:423}, \code{py:453}; it is not
a coordinate of $ζ$, so the argument is unaffected.) Neither $α$ nor $β$ is a coordinate of
$ζ$. If the zerocheck state $Z$ (every constraint extension, with the coordinates of $ζ$ drawn so
far substituted, is identically zero) is false, some $\tilde{C}_{t,i}$ has a nonzero restriction; $Z$
becomes true only if that one becomes zero: probability at most $1/|\E|$ on $x_i$ for
$i + 2 < τ_{\max}$, zero on $x_i$ for $i + 2 ≥ τ_{\max}$ (no table has that coordinate), and at
most $2/|\E|$ on the pair $(\mathit{low}, \mathit{high})$. \textbf{In total $τ_{\max}/|\E|$, once
for all constraints}: no union over the constraints is taken. In the bus phase's state
function, a conjunction $G ∧ Z$, the error of a challenge is $\max(ε_G, ε_Z)$, not the sum; on the
rounds of the last layer $ε_G = 5/|\E| ≥ 1/|\E|$, on the last pair $ε_G = 2/|\E| ≥ 2/|\E|$.
\textbf{The recycled point adds nothing to \lean{busError}}, per challenge or in the sum, when the
state function is the conjunction. Verdict on the blueprint: an upper bound, so not wrong; loose
by the factor $B$ and by the addition; attributed to $(α, β)$, which do not carry it; stated
three ways. Under the spine it can only be the bus phase's: a phase's state function charges that
phase's challenges, so acceptance test 6's ``the state function of Layer 7'' cannot be meant
literally.

\fhead{Proposed change}
\begin{change}{one charge, in the bus phase's state function}
\asis (\code{bp:333}) ``… is charged in the bus phase, coordinate by coordinate as $ζ$ is drawn,
$1/|\E|$ per coordinate per constraint …''
\asproposed ``… is carried by the bus phase's state function as one more conjunct, `every
constraint extension, restricted to the coordinates of $ζ$ drawn so far, is identically zero',
which a coordinate makes true with probability at most $1/|\E|$, whatever the number of
constraints. The coordinates of $ζ$ are the challenges of the last GKR layer (its rounds, then its
combination pair), whose own errors are at least that, so \lean{busError} has no term for it.''
Layer 7 (\code{bp:1000-1003}): replace ``it is charged to the $(α, β)$ and GKR challenges of
Layer 6, where $ζ$ is drawn, as an extra $τ_{\max}/|\E|$ per constraint'' by ``it is the bus
phase's (Layer 6), and costs $τ_{\max}/|\E|$ at most, once''. Acceptance test 6
(\code{bp:1283}): ``The state function of Layer 6 carries the zerocheck.''
\end{change}
\end{finding}

%------------------------------------------------------------------------------------------------
\begin{finding}{minor}{f:table-pub-layer7-sketch}{Layer 7's sketch and the tracker's signature predate the spine}
\fhead{Evidence} Layer 7's sketch against the spine as built (dossier section 4, C.2):

{\footnotesize\renewcommand{\arraystretch}{1.15}
\begin{tabular}{@{}>{\raggedright\arraybackslash}p{0.2\linewidth}>{\raggedright\arraybackslash}p{0.37\linewidth}>{\raggedright\arraybackslash}p{0.33\linewidth}@{}}
\toprule
 & Layer 7 (\code{bp:986-995}) with Layer 6's \lean{BusOut} (\code{bp:960-964}) & The spine as built \\
\midrule
Input statement & \lean{BusOut} = \code{ζ : Fin μ_bus → E}, \code{rem : Fin 3 → E},
  \code{pool : List Claim}, $α$, $β$ & \lean{I.Stmt × BusOut I}, \lean{BusOut I} =
  \lean{linear}, \lean{columns} \\
The public input & absent & threaded through every seam \\
The point $ζ$ & a field & absent; one point per term \\
The bus forms & recomputed by the table verifier from $α, β, ζ$ and the leaf layout & data: terms
  with weights in $\E$ \\
The zerocheck claims & implicit (the instance's constraints at $ζ$) & data: one linear claim
  each, if the bus phase emits them \\
The three values & \lean{rem} & the \lean{value} of three linear claims \\
Output & \code{BusOut × ColumnClaims} (keeps $ζ$, \lean{rem}, $α$, $β$) &
  \lean{I.Stmt × TableOut I} (column claims only) \\
\lean{tableSumcheck_relOut_}\allowbreak\lean{implies_constraints} & \code{(∀ j C, C̃_j(ζ_{<τ_j}) = 0 → …)} \code{-- see below}
  & no counterpart; the statement is a placeholder \\
\bottomrule
\end{tabular}}

The tracker's hole comment has a third signature,
\begin{leancode}
tableSumcheck I (S : Sumcheck.Def) : Phase.Def I (BusOut I) (TableOut I) (tableSpec I) (Seam.bus I) (Seam.table I)
\end{leancode}
with \lean{tableSummand I ζ ξ α β rem}: the spine has no \lean{tableSpec} and \lean{Phase.Def}
takes no relation. \lean{tableSumcheck_relOut_implies_constraints} (\code{bp:991-992}) has no
statement.

\emph{Which to keep.} Neither as it stands. Layer 7's needs $α, β$ and the selectors of the leaf
stack, which the abstract instance does not carry. The spine's is too general
(\cref{f:table-pub-bus-seam}). leanVM's own seam is between the two, the Rust's \code{BusVerify}
(\code{leaf.rs:849-860}):
\begin{srccode}
pub struct BusVerify {
    pub claims: Vec<ColumnClaim>,
    pub bytecode_claims: Vec<BytecodeClaim>,
    pub count_root: F192,
    /// The GKR point ζ, reused as the table sumcheck's eq point.
    pub point: Vec<F192>,
    /// `forms[side][table]`, for the zerocheck to settle.
    pub forms: [Vec<BusForm>; 3],
    /// Per side, what the tables' blocks owe its leaf claim: `Ṽ₀(ζ)` less the
    /// framework blocks' decomposition. Derived here, pinned by the batch's target.
    pub totals: [F192; 3],
}
\end{srccode}
\src{crates/lean\_vm/src/leaf.rs:849-860}
one point, the forms as data per side and table, one value per side, the column claims. Keep from
the spine: the public input threaded, the output reduced to column claims, $α, β$ and the leaf
layout private to the bus phase. Take from leanVM: the point in the statement and the forms
indexed by side and table (the proposed change of \cref{f:table-pub-bus-seam}).

\fhead{Proposed change} Rewrite the sketch on the spine's names,
\lean{def tableSumcheck (I) : Phase.Def I (I.Stmt × BusOut I) (I.Stmt × TableOut I)},
\lean{tableSumcheckComplete}, \lean{tableSumcheckSecurity}, with
\lean{tableSummand I (s : BusOut I) (ξ : E)}; delete the placeholder (what it meant is the bus
phase's zerocheck conjunct).
\end{finding}

%------------------------------------------------------------------------------------------------
\begin{finding}{minor}{f:table-pub-limb-claims-flock}{Layer 9's sketch makes the eighteen limb claims an input of Flock}
\fhead{Evidence} (For whoever reviews Layer 9.) Dossier section 6, E.2; blueprint
\code{bp:1078-1081}, \code{reduction … (StmtIn := ColumnClaims) … (StmtOut := WeightedClaim)},
``\lean{relIn} … the eighteen column claims are true of \lean{limbColumns}, and aux''. In leanVM
Flock's reduction takes no claim and the limb claims are opened with the others: the Rust's Flock
verifier takes no claim (\code{cpu/mod.rs:765}, \code{verify_reduction(n_blocks, &mut vs)}), and
nothing between the table sumcheck and the opening treats these claims apart
(\code{cpu/mod.rs:788-789}, ``No downstream special-casing: it folds into the one opening like
every other point claim.''). The spine's Flock slot is right: it passes the column claims on.
\end{finding}

%------------------------------------------------------------------------------------------------
\begin{finding}{minor}{f:table-pub-status-wrong}{Two status findings are wrong at the pin}
\fhead{Evidence and proposed change}
\begin{itemize}
\item \emph{The round message.} The status file's finding on the table round message (F6): ``the
  table round polynomial is a cubic sent whole, four nodes, three wire scalars
  (\code{constraints.rs:187-194, 267})''. The wire carries coefficients, not values at nodes
  (\code{constraints.rs:191-194}; \code{transcript.rs:56-71}: ``Send one sumcheck round
  polynomial, as its COEFFICIENTS, constant first''). The phrase copies a stale comment
  (\code{constraints.rs:24}; \cref{f:table-pub-stale-comments}). Proposed: ``a cubic, three of its
  four coefficients on the wire ($c₀, c₂, c₃$)''.
\item \emph{The caps in the Python verifier.} The status file's finding on the Python caps (F9):
  ``the Python verifier omits the caps $\mathit{log\_mem} ∈ [16, 32]$, $τ_j ≤ 32$, the bytecode
  power-of-two bound and $τ_{\mathrm{BLAKE2S}} ≥ 3$''. At the pin it has them
  (\code{py:857-864}):
\begin{srccode}
    log_bytecode = log2_strict(len(bytecode)) - BUS_BITS
    require(
        16 <= log_memory <= 32
        and all(0 <= log_height <= 32 for log_height in table_log_heights)
        and table_log_heights[OP_BLAKE2S] >= 3
        and 0 <= log_bytecode <= 32,
        "invalid announced table sizes",
    )
\end{srccode}
\src{python-verifier/verifier.py:857-864}
  called from \code{verify_execution} at \code{:1378}; \code{log2_strict} requires a power of two
  (\code{:252-254}). Outside the two phases of this dossier; recorded because a later part
  enumerates checks. Proposed: retire that finding (F9).
\end{itemize}
\end{finding}

%------------------------------------------------------------------------------------------------
\begin{finding}{minor}{f:table-pub-tracker-stale}{The tracker's text for the public-input phase is stale}
\fhead{Evidence} The hole comment says of the phase ``the prover sends nothing'' and ``the two
scalars are Layer 12's encoding to read and reconcile''; the phase merged with the prover sending
the values (decision 15). Text to draft only; nothing was posted.
\end{finding}

%------------------------------------------------------------------------------------------------
\begin{finding}{note}{f:table-pub-spec-errors}{The specification's batching error is loose by one; it gives no error for the recycled point; it does not say which coefficient is dropped}
\fhead{Evidence} What is stated (dossier section 5, D.1):
\begin{itemize}
\item \emph{Batching by $ξ$.} The specification charges $(ν_{\mathrm{side}} + B)/|\E|$
  (\code{05:155}, citing Corollary 3.9, \code{cor:idtest}, \code{03:67}); the blueprint
  $(B + 2)/|\E|$ on $ξ$ (\code{bp:994}). The tight value is $(k−1)/|\E| = (B+2)/|\E|$ for the
  $k = B + 3$ claims of the deployed input (dossier section 5, D.2): the specification's
  $(ν_{\mathrm{side}} + B)$ is an upper bound one too large. (The citation check corrected the
  number of the corollary, which the dossier gave as 3.7.)
\item \emph{Each round.} $3/|\E|$, $3·τ_{\max}$ in all (\code{05:155}, citing Fact 3.10,
  \code{fact:sumcheck}, \code{03:79}); the blueprint $3/|\E|$ per round (\code{bp:994}). (The
  citation check corrected the number of the fact, which the dossier gave as 3.8, the number of
  the Schwartz--Zippel lemma.)
\item \emph{The recycled point.} None in the specification: \code{05:132} says it is sound ``as
  long as it results from uniform sampling occurring after the table columns are committed''
  (the blueprint's three charges: \cref{f:table-pub-zerocheck-escape}).
\item \emph{Public input.} $1/|\E|$ (\code{08:33}); the blueprint $1/|\E|$ on the one challenge
  (\code{bp:1027}).
\item The specification has no round-by-round statement for these phases: its only one is Theorem
  B.7 (\code{thm:rbr}), for the opening (\file{doc/leanvm/body/b-polynomial-commitment-scheme.tex:139-152}),
  with the remark that the round-by-round error ``is the maximum of these entries over all
  rounds''. (The citation check corrected the number of the theorem, which the dossier gave as
  B.2.)
\item \emph{Which coefficient is dropped.} The specification does not say which coefficient is
  omitted; §3's ``a zerocheck round on a degree-$d$ cofactor therefore costs $d$ field elements''
  (\code{03:110}) does not apply to the table sumcheck, which sends the cubic (\code{08:76}),
  because the waiting tables' line $Y·u$ is not a multiple of the $\mathrm{eq}$ factor
  (\code{constraints.rs:20-26}).
\end{itemize}
\end{finding}

%------------------------------------------------------------------------------------------------
\begin{finding}{note}{f:table-pub-stale-comments}{Stale comments in the Rust on the round message}
\fhead{Evidence} \code{constraints.rs:24} says the round polynomial ``is sent WHOLE, at four
nodes''; \code{:265-266} says ``\code{h(0)} is derived from the running claim rather than
transmitted''. The code sends three \emph{coefficients} and derives $c₁$; $h(0) = c₀$ is
transmitted (\code{transcript.rs:63-71}). The blueprint follows the code; but the status repeats
``four nodes'' (\cref{f:table-pub-status-wrong}). The citation check found two more stale comments
in \code{constraints.rs}: \code{:26} says the verifier checks ``\code{h(0) + h(1) = claim}, then
interpolate at the challenge'' (there is neither a check nor an interpolation: coefficients are
read and one is derived); \code{:257-260} says ``Each round arrives as the round polynomial
itself at \code{nd}'' (it arrives as three coefficients, \code{transcript.rs:289-302}).
\end{finding}

%------------------------------------------------------------------------------------------------
\begin{finding}{note}{f:table-pub-slot-not-pinned}{The spine's slot does not pin the protocol: a phase that reads the whole oracle fills it}
\fhead{Evidence} Dossier section 6, E.4, item 6: \emph{Faithfulness itself.} In the ideal oracle
model a verifier may query $q$ anywhere. A ``phase'' with no message, whose verifier reads $q$ on
the whole cube and decides every claim, fills the slot with both proofs and error zero. The master
theorems hold for it. Nothing in the spine's types separates it from leanVM's sumcheck: that
rests on the phase's definition being reviewed against the transcript
(\cref{sec:gen-table-pub-transcript}), and on Layer 12's \lean{verify_iff_compiled} with its
fixture.
\end{finding}

%------------------------------------------------------------------------------------------------
\begin{finding}{note}{f:table-pub-other-notes}{Other notes: the column order inside a table, and perfect completeness}
\fhead{Evidence}
\begin{itemize}
\item \emph{The order of the columns inside a table} fixes the order of the 104 values, the ties
  of the stack and the powers of $λ$. The protocol blueprint does not cite it; the leanISA
  blueprint does (\file{docs/roadmap/leanisa-blueprint.md:612-613}; the citation check corrected
  the lines, which the dossier gave as 611-612). Whether leanISA's tables have the Rust's order,
  the two witnessed columns $w, b$ of JUMP last, was not checked.
\item \emph{Perfect completeness without exceptional challenges} (acceptance test 20) holds of
  both phases: the verifier's weights are products of $1 + ζ_m + r_m$ and $r_m$, the derivation
  of $c₁$ is a sum; the only inverse in the prover is of the constant $g + g²$
  (\code{primitives/src/multilinear.rs:233-241}).
\end{itemize}
\end{finding}
